Linear concept arithmetic is often evaluated on clean counterfactual pairs, but that setup leaves open a harder question: what happens along the straight-line path between endpoints that are semantically incompatible? We study this question in text embedding space with a two-stage design built on a real embedding model, \embmodel. First, we run an LRH-style separable-control experiment on curated word-pair assets and recover a strong linear regime for several transformations, including \texttt{Ving-3pSg} and adjective inflection contrasts. Second, we compare matched congruent and incongruent interpolation paths on 225 \sugarcrepe triplets spanning object swaps, attribute swaps, and relation replacements. We evaluate each path with off-manifold deviation, dual-space path length, dual detour ratio, switch count, and graph detour ratio.

The aggregate result contradicts a simple ``incongruent means more curved'' hypothesis: pooled incongruent paths have lower mean off-manifold deviation than congruent controls by 0.0083 and shorter dual paths by 0.0462. The key finding is heterogeneity. Relation replacements behave like genuine interpolation failures, increasing off-manifold deviation by 0.0154 and dual path length by 0.0785 relative to their matched paraphrase controls, while object and attribute swaps often look easier because the negative caption remains lexically close to the source. These results suggest that interpolation instability tracks structural incompatibility, especially disrupted relations, rather than contradiction alone. This narrows when linear traversal is a meaningful tool for representation analysis and latent steering.
